\subsection{Actual normalization of an original line-preimage component}
\label{subsec:normalized-line-components}

Let $F$ be a field, let $g\in F[X_i:i\in\sigma]$, and retain the actual
finite ambient algebra map $\varphi:F[X_i][1/g]\to B$. A direction with
$b_j\ne0$ defines the original transverse line equations. For a genuine
minimal prime $P$ above their original image in $B$, set
$D=B/P$ and $C=\overline D\subset\operatorname{Frac}(D)$, the literal
integral closure. This chapter constructs the normalization of this one
actual component; it does not replace the original full preimage by its
reduction, choose a component certificate, or assemble all components.

\begin{theorem}[The actual component normalization and its genuine free line action]
\label{algebra:normalized-line-component-ring}
\lean{EnvelopeProperty.AffineLine.componentQuotient_isDomain,
  EnvelopeProperty.AffineLine.componentQuotient_finiteType,
  EnvelopeProperty.AffineLine.normalizedComponentLineHom,
  EnvelopeProperty.AffineLine.normalizedComponentLineHom_principalLineHom,
  EnvelopeProperty.AffineLine.componentNormalization_inclusion_finite,
  EnvelopeProperty.AffineLine.normalizedComponentLineHom_finite,
  EnvelopeProperty.AffineLine.normalizedComponentLineHom_injective,
  EnvelopeProperty.AffineLine.componentNormalization_isDedekindDomain,
  EnvelopeProperty.AffineLine.openLineRing_isPrincipalIdealRing,
  EnvelopeProperty.AffineLine.normalizedComponentLineHom_free,
  EnvelopeProperty.AffineLine.normalizedComponentLineHom_flat}
\leanok
\uses{algebra:line-component-whole-map,algebra:line-component-domination,algebra:principal-open-line-ring,algebra:finite-free-curve-normalization}
The original minimal-prime hypothesis makes $D$ a domain. If $\sigma$ is
finite and $\varphi$ is finite, then $D$ is of finite type over $F$.
The actual component homomorphism
\[
 A=F[T][1/g(a+Tb)]\longrightarrow D\longrightarrow C
\]
retains every original localized ambient element. If $F$ is perfect, the
normalization inclusion and this whole line homomorphism are finite.
If $g(a+Tb)\ne0$, then $C$ is Dedekind. If also $B$ is a domain,
$g\ne0$ and $\varphi$ is injective, the line homomorphism is injective,
and its exact induced $A$-module structure on $C$ is free and flat.
No separability assumption is imposed.
\end{theorem}
\begin{proof}
A minimal prime is prime, hence its quotient is a domain. Give $B$ the
actual algebra structure induced by $\varphi$. Its scalar tower over $F$
is proved by the entire scalar identity of this algebra homomorphism.
The localization of the finite-variable polynomial ring is of finite type
over $F$; finiteness of $B$ over that ring then gives finite type of $B$
and of its actual quotient $D$.

Compose the already constructed whole component-line algebra homomorphism
with the literal inclusion $D\to C$. Its value on every original localized
ambient element is the inclusion of the original quotient value, by the
whole component factorization. Perfect-field finite normalization gives
finiteness of $D\to C$, and composition gives finiteness of $A\to C$.
Actual fraction-field inclusion is injective. The earlier genuine
component-domination theorem makes $A\to D$ injective under the stated
domain, nonvanishing and finite-injective ambient hypotheses, so the
composition remains injective.

The nonempty open-line ring is a Dedekind domain. Integral dimension
descent together with true normality and finite Noetherian normalization
makes $C$ Dedekind. Every ideal in $F[T][1/p]$ is the extension of its
contraction to $F[T]$; that contraction is principal and an ideal image
remains principal. Thus the actual open-line ring is a principal ideal
ring, even for the empty-open ring. For the nonempty line, the exact
injective action on $C$ gives torsion-freeness. The genuine finite
torsion-free principal-domain module theorem gives freeness, and free
modules are flat. The action in the freeness statement is explicitly the
one induced by the entire constructed normalized line homomorphism.
\end{proof}

\begin{theorem}[Entire normalized component scheme triangles]
\label{scheme:normalized-line-component-map}
\lean{EnvelopeProperty.AffineLine.normalizedComponentLineSchemeMap,
  EnvelopeProperty.AffineLine.normalizedComponentLineSchemeMap_factor,
  EnvelopeProperty.AffineLine.normalizedComponentLineSchemeMap_square,
  EnvelopeProperty.AffineLine.normalizedComponentLineSchemeMap_over_field,
  EnvelopeProperty.AffineLine.normalizedComponentLineSchemeMap_isFinite,
  EnvelopeProperty.AffineLine.normalizedComponentLineSchemeMap_flat,
  EnvelopeProperty.AffineLine.normalizedComponentLineSchemeMap_surjective}
\leanok
\uses{algebra:normalized-line-component-ring,algebra:line-component-scheme-domination,scheme:actual-finite-flat-curve-normalization}
The literal $\operatorname{Spec}C\to\operatorname{Spec}A$ is the actual
normalization projection followed by the original component-line scheme
map. Its whole ambient square and whole original $F$-structure triangle
commute. Under the corresponding ring hypotheses it is finite and flat.
The original domain, nonvanishing and finite-injective ambient hypotheses
make it surjective on genuine scheme points; this surjectivity itself
does not require perfection of $F$.
\end{theorem}
\begin{proof}
Apply the actual contravariant Spec functor to the entire normalization
and component ring composition. This proves the full factorization, not
merely its coordinate values. Compose the earlier original component
square with the literal normalization projection to obtain the complete
ambient square. The normalized algebra homomorphism's equality on every
$F$ scalar gives its entire structure triangle.

True finite and flat ring maps give finite and flat Spec maps. The
literal normalization projection is scheme-surjective by actual integral
lying over and its fraction-field injection. The original component-line
map is scheme-surjective by the genuine component-domination producer.
Their composition is therefore scheme-surjective, retaining the original
scheme map and all its triangles.
\end{proof}

\begin{theorem}[Native normalized line maps for every prescribed field valuation]
\label{scheme:normalized-line-native-map}
\lean{EnvelopeProperty.AffineLine.normalizedComponentLineNativeMap,
  EnvelopeProperty.AffineLine.continuous_normalizedComponentLineNativeMap,
  EnvelopeProperty.AffineLine.isProperMap_normalizedComponentLineNativeMap,
  EnvelopeProperty.AffineLine.normalizedComponentLineNativeMap_surjective}
\leanok
\uses{scheme:normalized-line-component-map,scheme:relative-integral-native-proper-map,scheme:relative-integral-native-image}
For any real-valued field valuation $\nu$ on $F$, the entire normalized
line morphism induces a continuous fixed-$\nu$ native point map. Under
finite-variable, perfect-field and finite ambient hypotheses it is
proper; adding the original domain, nonvanishing and injectivity
hypotheses makes it surjective. No bound on all elements of $F$, normed
field structure or completeness hypothesis is required.
\end{theorem}
\begin{proof}
The genuine entire $F$-structure triangle defines the actual relative
native point map and preserves every prescribed scalar value. Continuity
is the genuine native functoriality theorem. Finiteness of the actual
scheme map implies integrality, so the checked arbitrary prescribed-base
integral native properness theorem applies. Its scheme-surjectivity and
the exact common-base integral native image theorem give whole native
point lifts. There is no arbitrary different-prescribed-base extension
assertion in this argument.
\end{proof}

\begin{theorem}[Original ambient ideals through component normalization]
\label{scheme:normalized-component-original-ideals}
\lean{EnvelopeProperty.AffineLine.normalizedComponentInclusionSchemeMap,
  EnvelopeProperty.AffineLine.normalizedComponentInclusionSchemeMap_eq_specMap,
  EnvelopeProperty.AffineLine.normalizedComponentInclusion_idealSheaf_comap,
  EnvelopeProperty.AffineLine.normalizedComponentInclusion_stalkIdeal_comap}
\leanok
\uses{scheme:normalization-original-ideal-pullback,scheme:genuine-ideal-pullback-stalk}
The actual composite $\operatorname{Spec}C\to\operatorname{Spec}(B/P)
\to\operatorname{Spec}B$ pulls the original affine ideal datum of $J$
back to the full image $(J(B/P))C$. At every actual normalized point,
the pulled-back stalk ideal is the image of the original ideal stalk
under the entire actual stalk map. These identities require no
finite-generation or perfect-field hypothesis.
\end{theorem}
\begin{proof}
The genuine Spec composition identifies the composite with the actual
homomorphism $B\to B/P\to C$. Actual geometric ideal pullback on the
affine top open is the entire global-section ideal image. The genuine
$\Gamma\operatorname{Spec}$ naturality identity identifies this full
section map with the original ring composition, and ideal images compose.
Affine ideal sheaf extensionality gives the entire ideal datum equality.
The already proved arbitrary-scheme ideal pullback stalk theorem gives
the full stalk-map image at every actual normalized point. The original
$J$, including its powers and nilpotents, is not replaced by a radical.
These are geometric ideal identities, not ideal norm identities over
all coefficients of a possibly unbounded valued field.
\end{proof}

\begin{theorem}[Whole native lifts above the original component]
\label{scheme:normalization-component-native-lifts}
\lean{EnvelopeProperty.AffineLine.componentNormalizationNativeMap,
  EnvelopeProperty.AffineLine.componentNormalizationNativeMap_surjective,
  EnvelopeProperty.AffineLine.isProperMap_componentNormalizationNativeMap}
\leanok
\uses{scheme:actual-finite-flat-curve-normalization,scheme:relative-integral-native-proper-map,scheme:relative-integral-native-image}
For any prescribed field valuation, the actual normalization projection
induces a proper whole native component map. For the original genuine
minimal-prime quotient it is surjective, preserving every original
component value and every prescribed field scalar. These native
projection statements do not require finite type, perfection, or finite
module normalization.
\end{theorem}
\begin{proof}
The literal integral-closure inclusion is integral regardless of
finiteness. Its scalar-tower algebra homomorphism gives the true common
$F$-base triangle, so the general integral native properness theorem
applies. For the actual prime quotient, the fraction-field inclusion is
injective and integral lying over makes the normalization projection
scheme-surjective. The genuine exact common-base native integral image
theorem supplies whole prescribed-base lifts. A chosen branch above a
given scheme support is not asserted.
\end{proof}

\begin{theorem}[Given rational original points lift with every original value]
\label{scheme:normalized-component-rational-native-lift}
\lean{EnvelopeProperty.AffineLine.exists_componentNormalization_rationalEval,
  EnvelopeProperty.AffineLine.exists_componentNormalization_rationalNativeLift}
\leanok
\uses{algebra:normalized-line-component-ring,spectrum:integral-rational-evaluator-existence,scheme:affine-valued-point-homeomorphism}
Suppose $F$ is algebraically closed and a given actual evaluator
$e:B\to_F F$ kills the original component prime $P$. It extends to an
actual evaluator $d:C\to_F F$ satisfying
$d(\iota(\overline r))=e(r)$ for every original $r\in B$.
For any real-valued valuation $\nu$ of $F$, there is an actual
fixed-$\nu$ normalized native point whose whole image in
$\operatorname{Spec}B$ is the original rational point evaluated by
$\nu\circ e$, and whose entire affine valuation is $\nu\circ d$.
No finite-variable, finite-module, norm or completeness hypothesis is
required for this rational lift.
\end{theorem}
\begin{proof}
The exact kernel condition defines a genuine quotient evaluator
$B/P\to_F F$. The normalization inclusion is integral and injective into
the original component fraction field. Apply the checked actual rational
evaluator extension theorem over an algebraically closed field. Its
whole algebra-homomorphism identity gives the claimed equality on every
original $B$ element.

Compose $d$ with $\nu$ to obtain a genuine real-valued seminorm on $C$.
Its equality on all $F$ scalars gives an actual relative spectrum point,
and the genuine affine native homeomorphism gives a whole native point.
Apply actual affine naturality to its entire original component inclusion.
The evaluator equality on every original $B$ element and affine
homeomorphism injectivity identify its image with the original whole
native rational point. Its entire $C$ seminorm is $\nu\circ d$ by the
actual inverse affine homeomorphism law.

This constructs one actual normalized component and rational lifts, not
the normalization of the full reduced preimage, a smoothness theorem,
an arbitrary proper curve model, or an analytic potential comparison.
For nontrivial $\nu$, whole-field ideal norms must still not be substituted
for genuine bounded-model center-stalk functions. Those model and
potential constructions remain separate paper boundaries.
\end{proof}
